%!TEX TS-program = pdflatex
%!TEX TS-options = -shell-escape
% % % % %   Die folgenden Zeilen müssen ihre Zeilennummern 4 und 5 behalten !!!    % % % % %
\newcommand{\printpraesenzlsg}{false}
\newcommand{\printloesungen}{true}
\newcommand{\printbewertungen}{false}
% % % % %   \newcommand{\printloesungen}{false}                                    % % % % %
\newcommand{\blattnummer}{6}
%\newcommand{\abgabetermin}{\textcolor{red}{bis 11.04.2022, 09:00 Uhr}}
\documentclass[%
  paper=a4,
  fontsize=10pt,
  ngerman
  ]{scrartcl}

% Basics für Codierung und Sprache
% ===========================================================
  \usepackage{scrtime}
  \usepackage{etex}
  \usepackage{shellesc}
  \usepackage[final]{graphicx}
  \usepackage[utf8]{inputenc}
  \usepackage{babel}
  \usepackage[german=quotes]{csquotes}
  \usepackage{datetime2}
  \usepackage{ifthen}
  \usepackage{dingbat}
% ===========================================================

%% Fonts und Typographie
%% ===========================================================
  \usepackage{charter}
  \usepackage{nimbusmononarrow}
  \usepackage[babel=true,final,tracking=smallcaps]{microtype}
  \DisableLigatures{encoding = T1, family = tt* }
  \usepackage{ellipsis}
% ===========================================================

% Farben
% ===========================================================
  \usepackage[x11names,table]{xcolor}
  \definecolor{fbblau}{HTML}{3078AB}
  \definecolor{blackberry}{rgb}{0.53, 0.0, 0.25}
% ===========================================================

% Mathe-Pakete und -Einstellungen
% ===========================================================
  \usepackage{mathtools}
  \usepackage{amssymb}
  \usepackage[bigdelims]{newtxmath}
  \allowdisplaybreaks
  \usepackage{bm}
  \usepackage{wasysym}
  \usepackage{bussproofs}
  % skalierbarer Beweisbaum: \begin{scprooftree}{0.8} ... \end{scprooftree}
  \newenvironment{scprooftree}[1]%
    {\gdef\scalefactor{#1}\center\proofSkipAmount \leavevmode}%
    {\scalebox{\scalefactor}{\DisplayProof}\proofSkipAmount \endcenter}
% ===========================================================

% TikZ
% ===========================================================
  \usepackage{tikz}
  \usepackage{tikz-cd}
  \usepackage{pgf}
  \usetikzlibrary{matrix}
  \usetikzlibrary{positioning}
  \usetikzlibrary{arrows}
  \usetikzlibrary{shapes}
  \tikzset{>=Latex}
  \usepackage[mode=image]{standalone}
% ===========================================================

% Seitenlayout, Kopf-/Fußzeile
% ===========================================================
  \usepackage{scrlayer-scrpage}
  \usepackage[top=3cm, bottom=2.5cm, left=2.5cm, right=2cm]{geometry}  
  \pagestyle{scrheadings}
  \clearscrheadfoot 
  \setheadsepline{0.4pt}
  \setfootsepline{0.4pt}
  \setkomafont{pagehead}{\normalfont\footnotesize}
  \setkomafont{pagefoot}{\normalfont\footnotesize}
  \lohead[]{Klausur zur Vorlesung \textit{Mathematik I - Theoretische Grundlagen der Informatik} -- Wintersemester 2022/2023}
  \rohead[]{%
    \ifthenelse{\boolean{loesungen}}{\textbf{Lösungen zu Blatt \blattnummer}}{Blatt \blattnummer}}
  %\rofoot[Version: \DTMnow]{%
  %  \ifthenelse{\boolean{loesungen}}{Version: \DTMnow}{}}
  %\lofoot[\jobname.tex]{%
  %  \ifthenelse{\boolean{loesungen}}{\jobname.tex}{}}
  \cfoot[\thepage]{\thepage}
  \raggedbottom
  \usepackage{setspace}          
  \linespread{1.2}
  \setlength{\parindent}{0pt}
  \setlength{\parskip}{0.5\baselineskip}
  \usepackage[all]{nowidow}
  \usepackage{lscape}
% ===========================================================

% Hyperref
% ===========================================================
  \usepackage[%
    hidelinks,
    pdfpagelabels,
    bookmarksopen=true,
    bookmarksnumbered=true,
    linkcolor=black,
    urlcolor=fbblau,
    plainpages=false,
    pagebackref,
    citecolor=black,
    hypertexnames=true,
    pdfborderstyle={/S/U},
    linkbordercolor=fbblau,
    colorlinks=false,
    backref=false]{hyperref}
  \hypersetup{final}
  \makeatletter
  \g@addto@macro{\UrlBreaks}{\UrlOrds}
  \makeatother
% ===========================================================

% Listen und Tabellen
% ===========================================================
  \usepackage{multicol}
  \usepackage{multirow}
  \usepackage[shortlabels]{enumitem}
  \setlist[enumerate]{font=\bfseries,itemsep=0.25\baselineskip}
  \setlist[itemize]{label=$\triangleright$,itemsep=-0.25\baselineskip}
  \usepackage{tabularx}
  \usepackage{array}
  \usepackage{listings}
  \lstset{breaklines=true,
    basicstyle=\ttfamily,
    frame=tblr,
    language=c,
    numbers=left}
  \usepackage{include/mips}

\definecolor{lightgreen}{rgb}{0,0.7,0}
\definecolor{gray}{rgb}{0.5,0.5,0.5}
\definecolor{mauve}{rgb}{0.58,0,0.82}

\lstdefinestyle{mips}{ %
  language=[mips]Assembler,       % the language of the code
  keywordstyle=\color{blue},          % keyword style
  keywordstyle=[2]\color{magenta},     % keyword style
  keywordstyle=[3]\color{red},          % keyword style
  commentstyle=\color{lightgreen},       % comment style
  stringstyle=\color{mauve},         % string literal style
  escapeinside={\%*}{*)},            % if you want to add a comment within your code
  morekeywords={*,...}               % if you want to add more keywords to the set
}

\newcommand{\mipsinline}[1]{\lstinline[style=mips,columns=fixed]{#1}}
\newcommand{\mipsregister}[1]{\mipsinline{\$#1}}
  \newcommand{\qed}{\hfill \ensuremath{\Box}}
  \newcommand{\cmd}[1]{\texttt{#1}}
% ===========================================================

% Sonstiges
% ===========================================================
  \usepackage[textsize=small,color=Red1!80!OrangeRed1!80]{todonotes}
  \usepackage{lipsum}
% ===========================================================

% Shortcuts and Math Commands
% ===========================================================
  \newcommand{\BB}{\mathbb{B}}      % Bool
  \newcommand{\CC}{\mathbb{C}}      % komplexe Zahlen
  \newcommand{\NN}{\mathbb{N}}      % nat. Zahlen
  \newcommand{\QQ}{\mathbb{Q}}      % rat. Zahlen
  \newcommand{\RR}{\mathbb{R}}      % reelle Zahlen
  \newcommand{\ZZ}{\mathbb{Z}}      % ganze Zahlen
  \newcommand{\oh}{\mathcal{O}}      % Landau-O
  \newcommand{\ol}[1]{\overline{#1}}    % überstreichen
  \newcommand{\wt}[1]{\widetilde{#1}}    % überschlängeln
  \newcommand{\wh}[1]{\widehat{#1}}    % überdachen
  \newcommand{\entspricht}{\mathrel{\widehat{=}}}
  
  \DeclarePairedDelimiter{\absolut}{\lvert}{\rvert}    % Betrag
  \DeclarePairedDelimiter{\ceiling}{\lceil}{\rceil}    % aufrunden
  \DeclarePairedDelimiter{\Floor}{\lfloor}{\rfloor}    % aufrunden
  \DeclarePairedDelimiter{\Norm}{\lVert}{\rVert}      % Norm
  \DeclarePairedDelimiter{\sprod}{\langle}{\rangle}    % spitze Klammern
  \DeclarePairedDelimiter{\enbrace}{(}{)}          % runde Klammern
  \DeclarePairedDelimiter{\benbrace}{\lbrack}{\rbrack}  % eckige Klammern
  \DeclarePairedDelimiter{\penbrace}{\{}{\}}        % geschweifte Klammern
  \newcommand{\Underbrace}[2]{{\underbrace{#1}_{#2}}}   % bessere Unterklammerungen
  
  % Kurzschreibweisen für Faule und Code-Vervollständigung
  \newcommand{\abs}[1]{\absolut*{#1}}
  \newcommand{\ceil}[1]{\ceiling*{#1}}
  \newcommand{\flo}[1]{\Floor*{#1}}
  \newcommand{\no}[1]{\Norm*{#1}}
  \newcommand{\sk}[1]{\sprod*{#1}}
  \newcommand{\enb}[1]{\enbrace*{#1}}
  \newcommand{\penb}[1]{\penbrace*{#1}}
  \newcommand{\benb}[1]{\benbrace*{#1}}
  \newcommand{\stack}[2]{\makebox[1cm][c]{$\stackrel{#1}{#2}$}}

% Übungszettel-Krams
% ===========================================================
  \usepackage[tikz]{mdframed}
  \newmdenv[%
    backgroundcolor = gray!20,
    linewidth=0pt,
    skipabove = 0.5cm,
    skipbelow = 0cm
  ]{notes}
  
  % Lösungsausgabe
  \newboolean{praesenzlsgn}
  \setboolean{praesenzlsgn}{\printpraesenzlsg}
  \newboolean{loesungen}
  \setboolean{loesungen}{\printloesungen}
  % Bewertungsschema-Ausgabe
  \newboolean{bewertungen}
  \setboolean{bewertungen}{\printbewertungen}

  \newcommand{\iforiginal}[1]{\ifthenelse{\boolean{praesenzlsgn}}{}{#1}}
  \newcommand{\ifpraesenzlsg}[1]{\ifthenelse{\boolean{praesenzlsgn}}{#1}{}}
  \newcommand{\ifloesung}[1]{\ifthenelse{\boolean{loesungen}}{#1}{}}
  \newcommand{\ifbewertung}[1]{\ifthenelse{\boolean{bewertungen}}{#1}{}}

  \newcommand{\umbruchoriginal}{\iforiginal{\newpage}}        % Seitenumbruch nur im Original
  \newcommand{\umbruchpraesenzlsg}{\ifpraesenzlsg{\newpage}}  % Seitenumbruch nur in Lösung d. Präsenzaufgabe
  \newcommand{\umbruchlsg}{\ifloesung{\newpage}}              % Seitenumbruch nur in Lösung
  \newcommand{\umbruchbewertung}{\ifbewertung{\newpage}}      % Seitenumbruch nur in Lösung mit Bewertungsschema
  \newcommand{\umbruchnurpraesenz}{\ifthenelse{\boolean{loesungen}}{}{\umbruchpraesenzlsg}}
  \newcommand{\umbruchnurlsg}{\ifthenelse{\boolean{bewertungen}}{}{\umbruchlsg}}

  \usepackage{comment}
  \newenvironment{praesenzlsg}%
    {\ifthenelse{\boolean{praesenzlsgn}}{\textbf{--- Präsenzaufgabe Anfang ---} \mbox{} \newline}{}}%
	{\ifthenelse{\boolean{praesenzlsgn}}{\mbox{} \newline \textbf{--- Präsenzaufgabe Ende ---}}{}}
  \newenvironment{loesung}%
    {\ifthenelse{\boolean{loesungen}}{\textbf{--- Lösung Anfang ---} \mbox{} \newline}{}}%
    {\ifthenelse{\boolean{loesungen}}{\mbox{} \newline \textbf{--- Lösung Ende ---} \ifthenelse{\boolean{bewertungen}}{}{\umbruchlsg}}{}}
  \newenvironment{bewertung}%
    {\ifthenelse{\boolean{bewertungen}}{\textbf{--- Bewertungsschema Anfang ---} \mbox{} \newline}{}}%
    {\ifthenelse{\boolean{bewertungen}}{\mbox{} \newline \textbf{--- Bewertungsschema Ende ---} \umbruchlsg}{}}

  \ifthenelse{\boolean{praesenzlsgn}}{}{\excludecomment{praesenzlsg}}
  \ifthenelse{\boolean{loesungen}}{}{\excludecomment{loesung}}
  \ifthenelse{\boolean{bewertungen}}{}{\excludecomment{bewertung}}

  % Aufgabennummerierung
  \newcounter{aufgabennummer} 
  \setcounter{aufgabennummer}{1}
  \newcommand{\aufgabe}[1]{\textbf{Aufgabe \blattnummer.\theaufgabennummer} \stepcounter{aufgabennummer} \hfill (#1 Punkte)}
  %\newcommand{\aufgabetitel}[2]{\iforiginal{\vspace{0.5cm}} \textbf{Aufgabe \blattnummer.\theaufgabennummer \ (#2)} \stepcounter{aufgabennummer} \hfill (#1 Punkte)}
  \newcommand{\aufgabetitel}[2]{\iforiginal{\vspace{0.5cm}} \textbf{Aufgabe \theaufgabennummer \ (#2)} \stepcounter{aufgabennummer} \hfill (#1 Punkte)}
  \newcommand{\praesenztitel}[1]{\iforiginal{\vspace{0.5cm}} \textbf{Aufgabe \blattnummer.\theaufgabennummer \ (#1)} \stepcounter{aufgabennummer} \hfill \smallpencil}

\newcommand{\danger}{\raisebox{0.5mm}{\fontencoding{U}\fontfamily{futs}\selectfont\char 66\relax}}
\newcommand{\ddanger}[1]{{\danger} \textbf{#1} {\danger}}

% Rechnerstrukturen-Krams
% ===========================================================
  \newcommand{\Band}{\ensuremath{\mathbin{\texttt{AND}}}}
  \newcommand{\Bor}{\ensuremath{\mathbin{\texttt{OR}}}}
  \newcommand{\Bnot}{\ensuremath{\mathbin{\texttt{NOT}}}}
  \newcommand{\Bnand}{\ensuremath{\mathbin{\texttt{NAND}}}}
  \newcommand{\Bnor}{\ensuremath{\mathbin{\texttt{NOR}}}}
  \newcommand{\Bxor}{\ensuremath{\mathbin{\texttt{XOR}}}}
  \newcommand{\Bxnor}{\ensuremath{\mathbin{\texttt{XNOR}}}}
  \newcommand{\Bone}{\ensuremath{\mathbin{\texttt{ONE}}}}
  \newcommand{\Bzero}{\ensuremath{\mathbin{\texttt{ZERO}}}}
  \newcommand{\Binary}[1]{\ensuremath{\texttt{#1}}\xspace}
  \newcommand{\ieee}[3]{\fbox{\ensuremath{\texttt{#1}}}\fbox{\ensuremath{\texttt{#2}}}\fbox{\ensuremath{\texttt{#3}}}}
  \DeclareMathOperator{\DNF}{DNF}
  \DeclareMathOperator{\KNF}{KNF}
  \DeclareMathOperator{\xor}{\oplus}
  \newcommand{\dc}{\star} %don't-care-Wert
 
\usepackage{bussproofs}
% Aufbau und Aufgaben 1--7 nach blatt05.tex; Aufgaben 8--9 ergänzen das Blatt.
% Zeile 5: true = mit Lösungen, false = nur Aufgaben.
% Korrekte laufende Kopfzeile nur für dieses Blatt (gemeinsame Vorlage unverändert).
\lohead[]{Übungen zur Vorlesung \textit{Theoretische Informatik} -- Wintersemester 2026}
\begin{document}
\iforiginal{\thispagestyle{scrplain}
\vspace*{-3cm}
\begin{minipage}[t][1.1cm][c]{4.5cm}
  \includegraphics[width=4cm]{include/hwr-logo.png}
\end{minipage}
\hfill
\begin{minipage}[t][1.5cm][c]{8cm}
  \begin{center}
  \begin{footnotesize}
    \textsf{Prof. Dr.-Ing. Sebastian Schlesinger} \\[-0.1cm]
    \textsf{Fachbereich 2 - Duales Studium Wirtschaft \& Technik}
    %\url{https://www.hwr-berlin.de}
  \end{footnotesize}
  \end{center}
\end{minipage}
\hfill
%\begin{minipage}[t][1.6cm][c]{3.2cm}
%  \includegraphics[width=2.8cm]{include/ESlogo2}
%\end{minipage}

\vspace*{-0.3cm}
\begin{center} 
  \hrulefill \\[0.1cm]
  {\large Übungsblatt } \blattnummer \\[0.15cm]
  {\huge \bfseries Theoretische Informatik} \\[0.10cm]
  {HWR Berlin, Wintersemester 2026} \\[-0.4cm]
  \begin{tabular}{lcr}
    \hspace{0.3\textwidth}   & \hspace{0.3\textwidth} & \hspace{0.3\textwidth} \\
    Prof. Dr.-Ing. Sebastian Schlesinger   %&                        & "Ubung zur Klausurvorbereitung    \\ 
    %Blatt \blattnummer     % \abgabetermin          \\ 
  \end{tabular} \\[0.1cm]
  \hrulefill
\end{center}}

\begin{notes}\small\raggedright
\textbf{Konventionen:} DFA sind vollständig anzugeben; bei NFA dürfen
Übergangsmengen leer sein. In Übergangstabellen markiert $\to$ den Startzustand
und $*$ einen Endzustand. Ein Automat kann durch sein Tupel und eine
Übergangstabelle oder durch einen eindeutig beschrifteten Zustandsgraphen angegeben werden.
Bei der Potenzmengenkonstruktion genügen die erreichbaren Zustandsmengen.
\end{notes}

\aufgabetitel{$2$}{Formale Sprache modellieren}\\

Gegeben ist die Sprache $L = \{ w \in \{a,b\}^* \mid |w|_a=1 \}$, also die Anzahl der $a$ in $w$ ist 1.
Entwerfen Sie einen DFA, der $L$ erkennt.

\begin{loesung}
Ein vollständiger DFA ist $M=(Q,\Sigma,\delta,q_0,F)$ mit
$Q=\{q_0,q_1,q_2\}$, $\Sigma=\{a,b\}$, Startzustand $q_0$ und
$F=\{q_1\}$. Seine Übergänge sind:
\begin{center}
\begin{tabular}{c|cc}
 & $a$ & $b$\\\hline
$\to q_0$ & $q_1$ & $q_0$\\
$*q_1$ & $q_2$ & $q_1$\\
$q_2$ & $q_2$ & $q_2$
\end{tabular}
\end{center}
Die Zustände speichern, wie viele $a$ bisher gelesen wurden:
\begin{itemize}
  \item $q_0$: noch kein $a$,
  \item $q_1$: genau ein $a$,
  \item $q_2$: mindestens zwei $a$ (Fangzustand).
\end{itemize}
Ein $b$ verändert die Anzahl der $a$ nicht und lässt daher den Zustand
unverändert. Beim ersten $a$ wechselt der Automat von $q_0$ nach $q_1$,
beim zweiten von $q_1$ nach $q_2$. Danach kann kein weiteres Zeichen die
Anzahl wieder verringern, also bleibt der Automat in $q_2$.
Damit wird genau dann akzeptiert, wenn das Wort genau ein $a$ enthält.
Beispielsweise werden $a$, $ab$ und $bbabb$ akzeptiert;
$\varepsilon$, $bbb$ und $aba$ werden verworfen.
\end{loesung}

\aufgabetitel{$3$}{Minimierung von DFA}\\
Gegeben ist folgender DFA $M = (Q, \Sigma, \delta, q_0, F)$ mit:
\begin{itemize}
    \item $Q = \{A, B, C, D, E\}$
    \item $\Sigma = \{0,1\}$
    \item $\delta(A,0)=B,\ \delta(A,1)=C$
    \item $\delta(B,0)=A,\ \delta(B,1)=D$
    \item $\delta(C,0)=E,\ \delta(C,1)=A$
    \item $\delta(D,0)=E,\ \delta(D,1)=B$
    \item $\delta(E,0)=E,\ \delta(E,1)=E$
    \item $q_0 = A$
    \item $F = \{A, D\}$
\end{itemize}
Minimieren Sie den DFA.

\begin{loesung}
Alle Zustände sind erreichbar: $A$ durch $\varepsilon$, $B$ durch $0$,
$C$ durch $1$, $D$ durch $01$ und $E$ durch $10$.
Wir wenden den Tabellenalgorithmus auf die ungeordneten Paare verschiedener
Zustände an. Eine Markierung bedeutet: Die beiden Zustände sind nicht äquivalent.

\textbf{Anfangsmarkierung:} Da $F=\{A,D\}$, markieren wir alle Paare,
bei denen genau ein Zustand akzeptierend ist:
\[
\{A,B\},\ \{A,C\},\ \{A,E\},\ \{B,D\},\ \{C,D\},\ \{D,E\}.
\]
Diese Paare sind bereits durch das leere Wort $\varepsilon$ unterscheidbar.

\textbf{Erste Runde:} Wir markieren weitere Paare, sobald ihre Folgezustände
für mindestens ein Eingabezeichen ein bereits markiertes Paar bilden:
\begin{center}
\begin{tabular}{c|c|c}
Neues markiertes Paar & Eingabezeichen & Bereits markiertes Folgepaar\\\hline
$\{B,C\}$ & $0$ & $\{A,E\}$\\
$\{B,E\}$ & $0$ & $\{A,E\}$\\
$\{C,E\}$ & $1$ & $\{A,E\}$
\end{tabular}
\end{center}
\textbf{Zweite Runde:} Nun wird auch $\{A,D\}$ markiert: Bei Eingabe von $0$
sind die Folgezustände $B$ und $E$, deren Paar in der ersten Runde markiert wurde.
Eine unterscheidende Fortsetzung für $A$ und $D$ ist daher $00$.

Die fertige Tabelle lautet; $\times_0$ steht für die Anfangsmarkierung,
$\times_1$ bzw.\ $\times_2$ für die Markierung in der ersten bzw.\ zweiten Runde:
\begin{center}
\begin{tabular}{c|cccc}
 & $A$ & $B$ & $C$ & $D$\\\hline
$B$ & $\times_0$ & & &\\
$C$ & $\times_0$ & $\times_1$ & &\\
$D$ & $\times_2$ & $\times_0$ & $\times_0$ &\\
$E$ & $\times_0$ & $\times_1$ & $\times_1$ & $\times_0$
\end{tabular}
\end{center}
Nur die Felder unterhalb der Diagonale stehen für Zustandspaare.
Dort bleibt kein Paar unmarkiert. Daher können keine zwei Zustände
verschmolzen werden: Der gegebene DFA ist bereits minimal.

\emph{Hinweis:} Die Runden sind hier so dargestellt, dass jeweils nur
Markierungen aus vorherigen Runden verwendet werden. Verwendet man neue
Markierungen sofort, kann sich die Reihenfolge ändern; das Endergebnis bleibt gleich.
\end{loesung}

\aufgabetitel{$3$}{Minimierung von DFA}\\
Gegeben ist folgender DFA $M = (Q, \Sigma, \delta, q_0, F)$ mit:
\begin{itemize}
    \item $Q = \{A, B, C, D\}$
    \item $\Sigma = \{0,1\}$
    \item $\delta(A,0)=C,\ \delta(A,1)=B$
    \item $\delta(B,0)=D,\ \delta(B,1)=A$
    \item $\delta(C,0)=A,\ \delta(C,1)=D$
    \item $\delta(D,0)=B,\ \delta(D,1)=C$
    \item $q_0 = A$
    \item $F = \{A,C\}$
\end{itemize}
Minimieren Sie den DFA.

\begin{loesung}

Alle Zustände sind erreichbar: $A$ durch $\varepsilon$, $B$ durch $1$,
$C$ durch $0$ und $D$ durch $01$.
Wir wenden den Tabellenalgorithmus auf die ungeordneten Paare verschiedener
Zustände an.

\textbf{Anfangsmarkierung:} Da $F=\{A,C\}$, markieren wir alle Paare,
bei denen genau ein Zustand akzeptierend ist:
\[
\{A,B\},\quad \{A,D\},\quad \{B,C\},\quad \{C,D\}.
\]
Unmarkiert bleiben zunächst $\{A,C\}$ und $\{B,D\}$.

\textbf{Prüfung der unmarkierten Paare:} Für beide Eingabezeichen betrachten
wir jeweils das Paar der Folgezustände:
\begin{center}
\begin{tabular}{c|cc}
Unmarkiertes Paar & Folgepaar bei $0$ & Folgepaar bei $1$\\\hline
$\{A,C\}$ & $\{A,C\}$ & $\{B,D\}$\\
$\{B,D\}$ & $\{B,D\}$ & $\{A,C\}$
\end{tabular}
\end{center}
Keines dieser Folgepaare ist markiert. Es kommt daher keine neue Markierung
hinzu; der Algorithmus endet. Die fertige Zustandspaartabelle lautet:
\begin{center}
\begin{tabular}{c|ccc}
 & $A$ & $B$ & $C$\\\hline
$B$ & $\times_0$ & &\\
$C$ & $-$ & $\times_0$ &\\
$D$ & $\times_0$ & $-$ & $\times_0$
\end{tabular}
\end{center}
Hier bezeichnet $\times_0$ eine Anfangsmarkierung und $-$ ein unmarkiertes
Paar. Nur die Felder unterhalb der Diagonale stehen für Zustandspaare.
Somit sind $A$ und $C$ äquivalent, ebenso $B$ und $D$.
Wir verschmelzen diese Paare zu $X=\{A,C\}$ und $Y=\{B,D\}$.
Der minimale DFA hat die Zustände $X$ und $Y$,
Startzustand $X$, Endzustandsmenge $\{X\}$ und folgende Übergänge:
\begin{center}
\begin{tabular}{c|cc}
 & $0$ & $1$\\\hline
$\to *X$ & $X$ & $Y$\\
$Y$ & $Y$ & $X$
\end{tabular}
\end{center}
Die Zustände sind durch $\varepsilon$ unterscheidbar, also ist der DFA minimal.
Er erkennt genau die Wörter mit einer geraden Anzahl von Einsen.
\end{loesung}

\aufgabetitel{$5$}{NFA und DFA für eine Sprache}\\
Konstruieren Sie einen NFA für die Sprache
$L = \{ w \in \{0,1\}^* \mid \text{ das dritte Zeichen von rechts ist } 1 \}$.

Wandeln Sie den NFA in einen äquivalenten DFA um.

\begin{loesung}

Ein NFA hat $Q=\{q_0,q_1,q_2,q_3\}$, Startzustand $q_0$,
$F=\{q_3\}$ und die Übergangstabelle
\begin{center}
\begin{tabular}{c|cc}
 & $0$ & $1$\\\hline
$\to q_0$ & $\{q_0\}$ & $\{q_0,q_1\}$\\
$q_1$ & $\{q_2\}$ & $\{q_2\}$\\
$q_2$ & $\{q_3\}$ & $\{q_3\}$\\
$*q_3$ & $\varnothing$ & $\varnothing$
\end{tabular}
\end{center}
Er rät mit dem Übergang $q_0\xrightarrow{1}q_1$ die drittletzte Position
und liest danach genau zwei weitere Zeichen.

Die Potenzmengenkonstruktion liefert acht erreichbare Zustände.
Wir schreiben die Namen der enthaltenen NFA-Zustände hintereinander:
So steht $q_0q_1$ für $\{q_0,q_1\}$ und $q_0q_1q_2$ für
$\{q_0,q_1,q_2\}$. Auch $q_0$ bezeichnet in der folgenden DFA-Tabelle
die Zustandsmenge $\{q_0\}$.
\begin{center}
\begin{tabular}{c|cc}
 & $0$ & $1$\\\hline
$\to q_0$ & $q_0$ & $q_0q_1$\\
$q_0q_1$ & $q_0q_2$ & $q_0q_1q_2$\\
$q_0q_2$ & $q_0q_3$ & $q_0q_1q_3$\\
$q_0q_1q_2$ & $q_0q_2q_3$ & $q_0q_1q_2q_3$\\
$*q_0q_3$ & $q_0$ & $q_0q_1$\\
$*q_0q_1q_3$ & $q_0q_2$ & $q_0q_1q_2$\\
$*q_0q_2q_3$ & $q_0q_3$ & $q_0q_1q_3$\\
$*q_0q_1q_2q_3$ & $q_0q_2q_3$ & $q_0q_1q_2q_3$
\end{tabular}
\end{center}
Für jeden Übergang vereinigen wir die Folgezustände aller enthaltenen
NFA-Zustände. Beispielsweise ergibt sich bei Eingabe von $1$ aus
$\{q_0,q_1\}$ die Menge $\{q_0,q_1\}\cup\{q_2\}=\{q_0,q_1,q_2\}$.

Startzustand ist $q_0$. Genau die Zustandsmengen, die $q_3$ enthalten,
sind akzeptierend. Die acht Tabellenzustände sind in der angegebenen
Reihenfolge durch $\varepsilon$, $1$, $10$, $11$, $100$, $101$, $110$ und $111$
erreichbar. Die Tabelle enthält für beide Eingabezeichen alle Übergänge;
weitere Zustandsmengen entstehen nicht.
Wörter mit weniger als drei Zeichen werden verworfen, weil dann $q_3$
noch nicht erreicht werden kann.
\end{loesung}

\aufgabetitel{$3$}{Pumping Lemma}\\
Zeigen Sie mit Hilfe des Pumping-Lemmas, dass die Sprache
$L = \{ a^n b^{2n} \mid n\geq 0\}$ nicht regulär ist.

\begin{loesung}
Angenommen, $L$ sei regulär. Dann gibt es eine Pumping-Länge $p\geq1$.
Wir wählen $w=a^pb^{2p}\in L$ mit $|w|=3p\geq p$.
Sei $w=xyz$ eine beliebige Zerlegung mit $|xy|\leq p$ und $|y|\geq1$.
Dann liegt $y$ vollständig im ersten $a$-Block, also $y=a^k$ mit $1\leq k\leq p$.
Für $i=0$ folgt
\[
xy^0z=a^{p-k}b^{2p}.
\]
Dieses Wort enthält $p-k$ Zeichen $a$, aber weiterhin $2p$ Zeichen $b$.
Für die Zugehörigkeit zu $L$ müssten es genau $2(p-k)$ Zeichen $b$ sein.
Wegen $k\geq1$ gilt jedoch $2(p-k)<2p$, also $xy^0z\notin L$.
Dies gilt für jede zulässige Zerlegung und widerspricht dem Pumping-Lemma.
Also ist $L$ nicht regulär. \qed
\end{loesung}

% \aufgabetitel{$5$}{Reguläre Ausdrücke}\\
% Gegeben ist der reguläre Ausdruck
% $r = (ab^*)^*a$

% Geben Sie einen NFA an und einen äquivalenten DFA, die beide die von $r$ beschriebene Sprache erkennen.

% \begin{loesung}

% Es gilt $L(r)=\{a\}\cup\{aua\mid u\in\{a,b\}^*\}$:
% Ein nichtleeres Wort aus $(ab^*)^*$ beginnt mit $a$, und jedes mit $a$
% beginnende Wort lässt sich vor jedem weiteren $a$ in Blöcke $ab^*$ zerlegen.
% Das abschließende $a$ des Ausdrucks kommt noch hinzu.

% Ein NFA ohne $\varepsilon$-Übergänge hat Startzustand $s$, Endzustandsmenge
% $\{f\}$ und folgende Übergänge:
% \begin{center}
% \begin{tabular}{c|cc}
%  & $a$ & $b$\\\hline
% $\to s$ & $\{t,f\}$ & $\varnothing$\\
% $t$ & $\{t,f\}$ & $\{t\}$\\
% $*f$ & $\varnothing$ & $\varnothing$
% \end{tabular}
% \end{center}
% Der Übergang nach $f$ rät das letzte $a$; in $t$ können weitere Blöcke entstehen.
% Der zugehörige vollständige DFA hat die erreichbaren Zustände
% $X=\{s\}$, $Y=\{t,f\}$, $Z=\{t\}$ und $\bot=\varnothing$:
% \begin{center}
% \begin{tabular}{c|cc}
%  & $a$ & $b$\\\hline
% $\to X$ & $Y$ & $\bot$\\
% $*Y$ & $Y$ & $Z$\\
% $Z$ & $Y$ & $Z$\\
% $\bot$ & $\bot$ & $\bot$
% \end{tabular}
% \end{center}
% \end{loesung}

\aufgabetitel{$2$}{Grammatik}\\
Gegeben ist die Grammatik $G$:
$S\rightarrow aSb \mid bSa \mid \varepsilon$

Geben Sie die Elemente von $L_4:=\{w\in L(G)\mid |w|\leq 4\}$ an.

Ist $L(G)$ regulär? Begründen Sie.

\begin{loesung}

Jeder nichtabschließende Schritt fügt genau zwei Zeichen hinzu. Damit ist
\[
L_4=\{\varepsilon,ab,ba,aabb,abab,baba,bbaa\}.
\]
Die Sprache ist nicht regulär. Angenommen, sie sei regulär mit Pumping-Länge $p\geq1$.
Das Wort $w=a^pb^p$ liegt in $L(G)$: Man wendet $p$-mal $S\to aSb$ und dann
$S\to\varepsilon$ an. Für jede zulässige Zerlegung $w=xyz$ mit $|xy|\leq p$
und $|y|\geq1$ ist $y=a^k$ mit $k\geq1$.
Durch Abpumpen entsteht $xz=a^{p-k}b^p$.
Jedes Wort aus $L(G)$ enthält jedoch gleich viele $a$ wie $b$, denn jede
nichtabschließende Regel erzeugt genau eines von beiden. Also gilt $xz\notin L(G)$,
ein Widerspruch zum Pumping-Lemma. \qed
\end{loesung}


\ifthenelse{\boolean{loesungen}}{}{\newpage}
\aufgabetitel{$14$}{Kombinationsaufgabe: Grammatik, NFA, DFA und Minimierung}\\
Gegeben sei die reguläre Grammatik $G=(V,\Sigma,P,S)$ mit
$V=\{S,A,B,C,D,E\}$, $\Sigma=\{a,b\}$ und den Produktionen
\[
\begin{aligned}
S&\to aS\mid bS\mid aA\mid bB,\\
A&\to aC\mid bD, & B&\to aD\mid bC,\\
C&\to bE, & D&\to bE, & E&\to\varepsilon.
\end{aligned}
\]
Wie auf Blatt 5 sind Regeln der Form $X\to\varepsilon$ hier zugelassen.
\begin{enumerate}[a)]
\item Konstruieren Sie unmittelbar aus den Produktionen einen äquivalenten NFA
mit einem Zustand pro Variable. Geben Sie alle Übergänge, Startzustand und
Endzustände an. Erklären Sie, wo Nichtdeterminismus auftritt. (3 Punkte)
\item Wandeln Sie diesen NFA mit der Potenzmengenkonstruktion in einen
vollständigen DFA um. Dokumentieren Sie alle erreichbaren Zustandsmengen
und ihre Übergänge. Fassen Sie noch keine äquivalenten Zustände zusammen. (5 Punkte)
\item Minimieren Sie den DFA aus b) mit dem Tabellenalgorithmus. Geben Sie die
Zustandspaartabelle mit den Markierungsrunden und die Übergangstabelle
des minimalen DFA an.
Begründen Sie, warum keine weitere Zusammenfassung möglich ist. (4 Punkte)
\item Beschreiben Sie $L(G)$ in Mengenschreibweise und geben Sie einen
regulären Ausdruck dafür an. Erklären Sie die Bedeutung der Zustände
des minimalen DFA. (2 Punkte)
\end{enumerate}

\begin{loesung}
\begin{enumerate}[a)]
\item Eine Produktion $X\to cY$ wird zum Übergang $X\xrightarrow{c}Y$;
wegen $E\to\varepsilon$ ist $E$ akzeptierend. Es ist
$N=(V,\Sigma,\delta_N,S,\{E\})$ mit
\begin{center}
\begin{tabular}{c|cc}
 & $a$ & $b$\\\hline
$\to S$ & $\{S,A\}$ & $\{S,B\}$\\
$A$ & $\{C\}$ & $\{D\}$\\
$B$ & $\{D\}$ & $\{C\}$\\
$C$ & $\varnothing$ & $\{E\}$\\
$D$ & $\varnothing$ & $\{E\}$\\
$*E$ & $\varnothing$ & $\varnothing$
\end{tabular}
\end{center}
In $S$ gibt es für jedes Zeichen zwei mögliche Folgezustände.
Der NFA rät dort den Beginn der letzten drei Zeichen.
Die $\varepsilon$-Produktion macht $E$ zum Endzustand; sie erfordert hier
keinen $\varepsilon$-Übergang.

\newpage
\item Für eine Zustandsmenge $T$ gilt
$\delta_D(T,c)=\bigcup_{q\in T}\delta_N(q,c)$.
Startmenge ist $\{S\}$; genau die Mengen mit $E$ sind akzeptierend.
Die vollständige Tabelle der erreichbaren Mengen lautet:
\begin{center}
\begin{tabular}{c|l|cc|l}
 & Zustandsmenge & $a$ & $b$ & erreichbar durch\\\hline
$\to q_0$ & $\{S\}$ & $q_1$ & $q_2$ & $\varepsilon$\\
$q_1$ & $\{S,A\}$ & $q_3$ & $q_4$ & $a$\\
$q_2$ & $\{S,B\}$ & $q_5$ & $q_6$ & $b$\\
$q_3$ & $\{S,A,C\}$ & $q_3$ & $q_7$ & $aa$\\
$q_4$ & $\{S,B,D\}$ & $q_5$ & $q_8$ & $ab$\\
$q_5$ & $\{S,A,D\}$ & $q_3$ & $q_7$ & $ba$\\
$q_6$ & $\{S,B,C\}$ & $q_5$ & $q_8$ & $bb$\\
$*q_7$ & $\{S,B,D,E\}$ & $q_5$ & $q_8$ & $aab$\\
$*q_8$ & $\{S,B,C,E\}$ & $q_5$ & $q_8$ & $abb$
\end{tabular}
\end{center}
Beispielsweise ist
\[
\delta_D(\{S,A,C\},b)
=\{S,B\}\cup\{D\}\cup\{E\}=\{S,B,D,E\}.
\]
Die Tabelle enthält alle erreichbaren Mengen und ist unter beiden Übergängen
abgeschlossen. Da $S$ immer enthalten bleibt, ist $\varnothing$ nicht erreichbar.

\newpage
\item Wir markieren die ungeordneten Paare verschiedener Zustände wie in
Aufgabe 2 und 3. Alle neun Zustände sind nach b) erreichbar.

\textbf{Anfangsmarkierung ($\times_0$):} Genau $q_7$ und $q_8$ sind
akzeptierend. Deshalb markieren wir die 14 Paare $\{q_i,q_j\}$ mit
$i\in\{0,\ldots,6\}$ und $j\in\{7,8\}$.

\textbf{Erste Runde ($\times_1$):} Wir markieren die zwölf Paare
$\{q_i,q_j\}$ mit $i\in\{0,1,2\}$ und $j\in\{3,4,5,6\}$.
Bei Eingabe von $b$ führt $q_i$ in einen nichtakzeptierenden Zustand
($q_2$, $q_4$ oder $q_6$), $q_j$ hingegen nach $q_7$ oder $q_8$.
Das jeweilige Folgepaar ist daher bereits mit $\times_0$ markiert.
Zum Beispiel führt $\{q_0,q_3\}$ bei $b$ zum markierten Paar $\{q_2,q_7\}$.

\textbf{Zweite Runde ($\times_2$):} Es kommen zwei Markierungen hinzu:
\begin{center}
\begin{tabular}{c|c|c}
Neues markiertes Paar & Eingabezeichen & Bereits markiertes Folgepaar\\\hline
$\{q_0,q_1\}$ & $a$ & $\{q_1,q_3\}$\\
$\{q_0,q_2\}$ & $a$ & $\{q_1,q_5\}$
\end{tabular}
\end{center}
Die fertige Tabelle lautet ($-$ bedeutet unmarkiert):
\begin{center}
\begin{tabular}{c|cccccccc}
 & $q_0$ & $q_1$ & $q_2$ & $q_3$ & $q_4$ & $q_5$ & $q_6$ & $q_7$\\\hline
$q_1$ & $\times_2$ &  &  &  &  &  &  & \\
$q_2$ & $\times_2$ & $-$ &  &  &  &  &  & \\
$q_3$ & $\times_1$ & $\times_1$ & $\times_1$ &  &  &  &  & \\
$q_4$ & $\times_1$ & $\times_1$ & $\times_1$ & $-$ &  &  &  & \\
$q_5$ & $\times_1$ & $\times_1$ & $\times_1$ & $-$ & $-$ &  &  & \\
$q_6$ & $\times_1$ & $\times_1$ & $\times_1$ & $-$ & $-$ & $-$ &  & \\
$q_7$ & $\times_0$ & $\times_0$ & $\times_0$ & $\times_0$ & $\times_0$ & $\times_0$ & $\times_0$ & \\
$q_8$ & $\times_0$ & $\times_0$ & $\times_0$ & $\times_0$ & $\times_0$ & $\times_0$ & $\times_0$ & $-$\\
\end{tabular}
\end{center}
Nur die Felder unterhalb der Diagonale stehen für Zustandspaare.
Es entstehen keine weiteren Markierungen:
\begin{itemize}
\item $\{q_1,q_2\}$ führt bei $a$ zu $\{q_3,q_5\}$ und bei $b$ zu
$\{q_4,q_6\}$; beide Paare bleiben unmarkiert.
\item Bei Paaren aus $q_3,q_4,q_5,q_6$ sind die Folgezustände entweder gleich
oder bilden bei $a$ das unmarkierte Paar $\{q_3,q_5\}$ bzw.\ bei $b$ das
unmarkierte Paar $\{q_7,q_8\}$.
\item $q_7$ und $q_8$ haben für jedes Zeichen denselben Folgezustand.
\end{itemize}
Wir verschmelzen die unmarkierten Paare zu den Zuständen
\[
K_0=\{q_0\},\quad K_1=\{q_1,q_2\},\quad
K_2=\{q_3,q_4,q_5,q_6\},\quad K_3=\{q_7,q_8\}.
\]
Startzustand ist $K_0$, einziger Endzustand ist $K_3$.
Die Übergangstabelle des minimalen DFA lautet:
\begin{center}
\begin{tabular}{c|cc}
 & $a$ & $b$\\\hline
$\to K_0$ & $K_1$ & $K_1$\\
$K_1$ & $K_2$ & $K_2$\\
$K_2$ & $K_2$ & $K_3$\\
$*K_3$ & $K_2$ & $K_3$
\end{tabular}
\end{center}
Alle vier Zustände sind erreichbar (z.\,B.\ durch $\varepsilon,a,aa,aab$).
$K_3$ ist von allen anderen durch $\varepsilon$ unterscheidbar.
$K_2$ unterscheidet sich von $K_0$ und $K_1$ durch die Fortsetzung $b$.
$K_0$ und $K_1$ unterscheiden sich durch $ab$.
Somit sind alle Paare unterscheidbar und der DFA ist minimal.

\newpage
\item Es gilt
\[
L(G)=\{w\in\{a,b\}^*\mid |w|\geq3\text{ und $w$ endet auf }b\},
\qquad r=(a\mid b)^*(a\mid b)(a\mid b)b.
\]
In $S$ entsteht ein beliebiges Präfix. Nach Verlassen von $S$ entstehen
genau drei Zeichen: Die ersten beiden sind beliebig, das dritte ist $b$.
Umgekehrt erhält man jedes solche Wort, indem man sein Präfix in $S$ erzeugt
und die letzten drei Zeichen über $A$ bzw.\ $B$ und $C$ bzw.\ $D$ ableitet.

$K_0$ steht für das leere Wort, $K_1$ für Wörter der Länge $1$,
$K_2$ für Wörter der Länge $2$ sowie längere Wörter mit letztem Zeichen $a$.
$K_3$ steht für Wörter der Länge mindestens $3$ mit letztem Zeichen $b$.

\end{enumerate}
\end{loesung}

% \aufgabetitel{$6$}{Pumping-Lemma: zwei Schranken}\\
% Betrachten Sie die Sprache
% \[
% L=\{a^n b^m\mid n,m\in\mathbb{N}_0,\ n<m\leq2n\}.
% \]
% \begin{enumerate}[a)]
% \item Zeigen Sie mit dem Pumping-Lemma für reguläre Sprachen, dass $L$ nicht
% regulär ist. Wählen Sie ein geeignetes Wort in Abhängigkeit von der
% Pumping-Länge $p$ und argumentieren Sie für \emph{jede} zulässige Zerlegung. (4 Punkte)
% \item Jemand wählt $w=a^pb^{p+1}$ und behauptet, Abpumpen ($i=0$) führe
% für jede zulässige Zerlegung aus $L$ heraus. Prüfen Sie diese Behauptung.
% Kann man mit demselben Wort und einem anderen Pumpfaktor den Beweis führen? (2 Punkte)
% \end{enumerate}

% \begin{loesung}
% \begin{enumerate}[a)]
% \item Angenommen, $L$ sei regulär. Dann existiert eine Pumping-Länge $p\geq1$.
% Wir wählen das Wort an der oberen Schranke:
% \[
% w=a^pb^{2p}\in L,\qquad |w|=3p\geq p.
% \]
% Sei $w=xyz$ eine beliebige Zerlegung mit $|xy|\leq p$ und $|y|\geq1$.
% Da die ersten $p$ Zeichen alle $a$ sind, ist $y=a^k$ mit $1\leq k\leq p$.
% Für $i=0$ ist
% \[
% xy^0z=a^{p-k}b^{2p}.
% \]
% Dieses Wort liegt nicht in $L$, denn die obere Schranke wäre
% $2p\leq2(p-k)$, während wegen $k\geq1$ tatsächlich $2p>2(p-k)$ gilt.
% Damit scheitert für jede zulässige Zerlegung mindestens ein Pumpfaktor.
% Das widerspricht dem Pumping-Lemma, also ist $L$ nicht regulär. \qed

% \item Die Behauptung zum Abpumpen ist falsch. Für $p=3$ ist etwa
% $w=a^3b^4$ und die Zerlegung $x=\varepsilon$, $y=a$, $z=a^2b^4$ zulässig.
% Nach Abpumpen bleibt $a^2b^4\in L$, denn $2<4\leq4$.
% Allgemein funktioniert dieses Gegenbeispiel mit $|y|=1$ für jedes $p\geq3$.
% Das widerlegt nur die vorgeschlagene Beweisführung, nicht die Nichtregularität.

% Mit $i=2$ gelingt der Beweis für dasselbe Wort: Für jede zulässige Zerlegung
% ist wieder $y=a^k$ mit $k\geq1$, und
% \[
% xy^2z=a^{p+k}b^{p+1}\notin L,
% \]
% weil die erforderliche strikte Ungleichung $p+k<p+1$ verletzt ist.
% $w=a^pb^{p+1}$ liegt für jedes $p\geq1$ in $L$, da $p<p+1\leq2p$.
% Somit liefert auch diese Wahl einen vollständigen Widerspruchsbeweis.
% \end{enumerate}
% \end{loesung}

\aufgabetitel{$5$}{Grammatik zu einer nichtregulären Sprache}\\
Gegeben ist die nichtreguläre Sprache
\[
L=\{a^n b^m c^n\mid n,m\in\mathbb{N}_0\}.
\]
Die Wörter bestehen aus einem $a$-Block, einem $b$-Block und einem $c$-Block.
Dabei gibt es gleich viele $a$ wie $c$; die Anzahl der $b$ ist unabhängig davon.
Einzelne Blöcke dürfen leer sein.
\begin{enumerate}[a)]
\item Geben Sie eine kontextfreie Grammatik $G$ mit $L(G)=L$ an. (2 Punkte)
\item Geben Sie mit Ihrer Grammatik Ableitungen für $aabbcc$, $bbb$ und
$\varepsilon$ an. (1 Punkt)
\item Begründen Sie, warum Ihre Grammatik genau die Wörter aus $L$ erzeugt.
Berücksichtigen Sie beide Inklusionen. (2 Punkte)
\end{enumerate}
Die Nichtregularität von $L$ darf als bekannt vorausgesetzt werden.

\begin{loesung}
\begin{enumerate}[a)]
\item Eine mögliche Grammatik ist $G=(\{S,B\},\{a,b,c\},P,S)$ mit
\[
P:\qquad S\to aSc\mid B,\qquad B\to bB\mid\varepsilon.
\]
Die Regel $S\to aSc$ erzeugt jeweils ein zusammengehöriges $a$ und $c$.
Nach dem Wechsel zu $B$ wird der mittlere $b$-Block erzeugt.
Da jede linke Seite eine einzelne Variable ist, ist die Grammatik kontextfrei.

\item Die gewünschten Ableitungen lauten:
\begin{align*}
S&\Rightarrow aSc\Rightarrow aaScc\Rightarrow aaBcc
  \Rightarrow aabBcc\Rightarrow aabbBcc\Rightarrow aabbcc,\\
S&\Rightarrow B\Rightarrow bB\Rightarrow bbB\Rightarrow bbbB
  \Rightarrow bbb,\\
S&\Rightarrow B\Rightarrow\varepsilon.
\end{align*}

\item \emph{$L(G)\subseteq L$:} Solange $S$ vorkommt, haben die Satzformen
 die Gestalt $a^nSc^n$. Nach $S\to B$ haben sie die Gestalt $a^nb^mBc^n$.
Die einzige Möglichkeit, die letzte Variable zu entfernen, ist $B\to\varepsilon$.
Jedes erzeugte Wort hat daher die Form $a^nb^mc^n$ mit $n,m\geq0$ und liegt in $L$.

\emph{$L\subseteq L(G)$:} Seien $n,m\geq0$ beliebig.
Wir wenden $n$-mal $S\to aSc$, dann einmal $S\to B$,
$m$-mal $B\to bB$ und schließlich $B\to\varepsilon$ an:
\[
S\Rightarrow^* a^nSc^n\Rightarrow a^nBc^n
 \Rightarrow^* a^nb^mBc^n\Rightarrow a^nb^mc^n.
\]
Dieses Schema gilt auch für $n=0$ oder $m=0$; die jeweilige Regel wird dann
keinmal angewendet. Also ist jedes Wort aus $L$ ableitbar und insgesamt $L(G)=L$. \qed
\end{enumerate}
\end{loesung}

\end{document}
